\begin{titlepage}
\thispagestyle{empty}
\begin{center}
    \vspace*{1.0cm}
    {\Large\bfseries\color{physCyan} เอกสารคำสอนและคู่มือปฏิบัติการรายวิชา 4011111}\\[0.5cm]
    {\Huge\bfseries\color{physRoyalBlue} ฟิสิกส์สำหรับจุลชีววิทยา}\\[0.4cm]
    {\LARGE\bfseries\color{physAmber} Physics for Microbiology}\\[0.4cm]
    {\large\bfseries\color{bioEmerald} หน่วยกิต 3(2-3-4)\quad หลักสูตรวิทยาศาสตรบัณฑิต สาขาวิชาจุลชีววิทยา}\\[1.6cm]

    \begin{tcolorbox}[
        enhanced,
        colback=boxBlueBg,
        colframe=physRoyalBlue,
        width=14.2cm,
        arc=2.5mm,
        boxrule=1.0pt,
        halign=center
    ]
        {\normalsize\bfseries ออกแบบตามแนวทางการจัดการศึกษาที่มุ่งเน้นผลลัพธ์ (Outcome-Based Education)}\\[4pt]
        {\small เชื่อมโยงผลลัพธ์การเรียนรู้ของหลักสูตร (PLO1--PLO6) สู่ผลลัพธ์การเรียนรู้รายวิชา (CLO1--CLO4)}\\[2pt]
        {\small พร้อมเทคนิคการคำนวณฉบับเข้าใจง่ายไร้แคลคูลัส และปฏิบัติการทดลอง 13 การทดลอง}
    \end{tcolorbox}

    \vspace{1.6cm}
    {\Large\bfseries ผู้ช่วยศาสตราจารย์ ดร.ชีวะ ทัศนา}\\[0.3cm]
    {\normalsize วท.บ. (ฟิสิกส์), วท.ม. (ฟิสิกส์ประยุกต์), ปร.ด. (ฟิสิกส์ประยุกต์)}\\[0.2cm]
    {\normalsize สาขาวิชาฟิสิกส์ คณะวิทยาศาสตร์และเทคโนโลยี}\\[0.2cm]
    {\normalsize มหาวิทยาลัยราชภัฏรำไพพรรณี}\\[1.8cm]

    \begin{tcolorbox}[
        enhanced,
        colback=softSmoke,
        colframe=bioEmerald,
        width=14.2cm,
        arc=2mm,
        halign=center
    ]
        {\small\bfseries เอกสารคำสอนประจำภาคเรียนที่ 2 ปีการศึกษา 2569}\\[2pt]
        {\footnotesize คณะวิทยาศาสตร์และเทคโนโลยี มหาวิทยาลัยราชภัฏรำไพพรรณี}\\[2pt]
        {\footnotesize พ.ศ. 2569}
    \end{tcolorbox}
\end{center}
\end{titlepage}
